\begin{abstract}
Multi-teacher on-policy distillation combines supervision from specialized teachers through one evolving student. We study how loss averaging, optimizer state, and vocabulary supervision shape this process using an aligned Qwen3-1.7B experiment with four domain teachers. Twelve capability evaluations show distinct early tradeoffs among three averaging rules: at update 50, global token averaging improves mathematics over domain response averaging by 4.6 percentage points, while domain token averaging reduces instruction-following accuracy by 2.0 points. Parameter measurements reveal a second distinction. After 100 sampled-gradient updates, 94.0\% of FP32 master coordinates have changed, compared with 3.41\% of BF16 model coordinates. At this checkpoint, saved Adam moments produce substantial movement even with a zero current gradient. Fixed-prefix comparisons further show that nearly identical writeback densities can conceal different update directions: sampled and full-vocabulary gradients have mean cosine 0.677, while their BF16 writeback steps have cosine 0.915. Together, these results connect capability allocation to the quantities actually averaged, accumulated, and stored during distillation, and identify input composition and optimizer history as essential variables in interpreting teacher-conditioned parameter geometry.
\end{abstract}
